\documentclass[a4paper]{article}

\usepackage[fontset=fandol]{ctex}
\usepackage{amsmath, amssymb}
\usepackage{graphicx}
\usepackage[margin=2.5cm]{geometry}
\usepackage[most]{tcolorbox}
\usepackage{etoolbox}
\usepackage{listings}
\usepackage{hyperref}
\usepackage{booktabs}
\usepackage{subcaption}
\usepackage{float}
\usepackage{tikz}
\IfFileExists{pgfplots.sty}{
    \usepackage{pgfplots}
    \pgfplotsset{compat=1.18}
}{}

\newtcolorbox{knowledgebox}[1]{
    enhanced, colback=blue!5!white, colframe=blue!75!black, colbacktitle=blue!75!black,
    coltitle=white, fonttitle=\bfseries, title=#1, attach boxed title to top left={yshift=-2mm, xshift=2mm},
    boxrule=1pt, sharp corners
}

\newtcolorbox{importantbox}[1]{
    enhanced, colback=yellow!10!white, colframe=yellow!80!black, colbacktitle=yellow!80!black,
    coltitle=black, fonttitle=\bfseries, title=#1, sharp corners
}

\newtcolorbox{warningbox}[1]{
    enhanced, colback=red!5!white, colframe=red!75!black, colbacktitle=red!75!black,
    coltitle=white, fonttitle=\bfseries, title=#1, sharp corners
}

\lstset{
    language=Python,
    basicstyle=\ttfamily\small,
    keywordstyle=\color{blue},
    stringstyle=\color{red!60!black},
    commentstyle=\color{green!60!black},
    breaklines=true,
    frame=single,
    numbers=left,
    numberstyle=\tiny\color{gray},
    captionpos=b,
    extendedchars=false
}

\newcommand{\notetitle}{CS224R 课程笔记：Lecture 16 RL for Robots}
\newcommand{\noteauthors}{Codex}
\newcommand{\notedate}{2026-04-13}
\newcommand{\videochannel}{Stanford Online}
\newcommand{\videopublishdate}{讲次日期 2025-05-23；YouTube 上传 2025-12-08}
\newcommand{\videoduration}{1:05:44}
\newcommand{\videourl}{https://www.youtube.com/watch?v=rbaWQQLrzl0}
\newcommand{\videocoverpath}{cover.jpg}

\begin{document}

\begin{titlepage}
\centering
{\Large 课程笔记\par}
\vspace{1.2cm}
{\huge\bfseries \notetitle\par}
\vspace{0.8cm}
{\large \noteauthors\par}
\vspace{0.3cm}
{\large \notedate\par}
\vspace{1.2cm}

\ifdefempty{\videocoverpath}{
{\small 请在生成笔记时填入视频封面图的本地路径。\par}
}{
\includegraphics[width=0.82\textwidth,height=0.45\textheight,keepaspectratio]{\videocoverpath}\par
}

\vfill
\begin{tcolorbox}[width=0.9\textwidth, colback=black!2!white, colframe=black!60, sharp corners]
\textbf{视频作者/频道}：\videochannel\par
\textbf{发布日期}：\videopublishdate\par
\textbf{视频时长}：\videoduration\par
\textbf{视频链接}：\href{\videourl}{\nolinkurl{\videourl}}
\end{tcolorbox}
\end{titlepage}

\tableofcontents
\newpage

\section{问题背景：为什么机器人还没有实现真正自主学习}

这讲围绕一个非常具体、但常被标准 RL 设定掩盖的问题展开：机器人当然可以在 rollout 期间根据当前策略自主执行动作，但一旦任务失败或偏离起点，究竟是谁把它放回可以重新尝试的位置？在仿真里，这一步被一句 \texttt{env.reset()} 掩盖；在真实世界里，这一步往往意味着人类重新摆放物体、扶正机器人、关上门、把毛巾展开，或者把移动平台重新开回起点。

\begin{figure}[H]
\centering
\includegraphics[width=0.95\textwidth]{pdf_pages/page-06.png}
\caption{标准强化学习默认存在环境重置接口；这一前提在真实机器人场景中往往不成立。}
\end{figure}

讲者把这一点表述得很直接：传统 RL 并不是完全自主的，只是把“自主”限定在一次 episode 内。机器人学里真正稀缺的不是动作执行能力，而是长时间、无人工干预的练习能力。对于导航、抓取、开门、整理物体等任务，能够反复自主尝试，比单次成功更接近长期部署时真正需要的能力。

\subsection{自主 RL 的定义}

本讲把 autonomous RL 定义为一种更严格的交互模式：环境在开始时初始化一次，之后智能体持续与世界交互，不依赖频繁环境重置。这个定义允许偶尔低频 reset，但关键思想是不再把 reset 当成免费操作。

\begin{figure}[H]
\centering
\includegraphics[width=0.82\textwidth]{pdf_pages/page-11.png}
\caption{自主 RL 的核心要求：初始化一次，然后在尽量没有环境重置的前提下持续学习。}
\end{figure}

这个定义把“训练时 autonomy”放到和“部署时 autonomy”同样重要的位置。讲者特别强调，机器人不仅要在任务执行时减少人类介入，也要在训练过程中减少人类介入，因为只有这样，真实世界数据规模才可能显著扩大，而数据规模扩大又可能带来更好的泛化和鲁棒性。

\subsection{评价指标也随之改变}

一旦去掉频繁 reset，评价标准也不再只有“最后学到的策略有多好”。本讲区分了两种目标：
\begin{itemize}
\item \textbf{deployed policy evaluation}：最终学到的策略在标准起点上有多强，适合关心最终性能的任务，例如机器人厨师。
\item \textbf{continuing policy evaluation}：智能体在整段生命期里累计拿到了多少奖励，适合关心持续产出的任务，例如火星车。
\end{itemize}

这两个目标都合理，但对应不同算法偏好。前者允许训练时为“最终好策略”忍受很多弯路；后者则要求学习过程本身就尽量少犯代价高昂的错误。

\begin{knowledgebox}{本讲的关键转向}
把真实世界 RL 看成“单次 episode 内优化动作”是不够的。更准确的视角是：智能体需要在很长的生命期中，一边探索、一边恢复、一边继续完成任务。
\end{knowledgebox}

\subsection{本章小结}

这一部分把问题重新摆正了：机器人 RL 的核心瓶颈并不只是策略优化，而是如何在没有人类频繁重置环境的前提下持续训练。自主 RL 的定义和评价方式都因此发生变化。

\section{长时程自主学习的难点}

\subsection{为什么简单地把 episode 拉长并不够}

一个自然想法是：既然 reset 代价高，那就把 horizon \(H\) 拉长，让 agent 少 reset。讲者指出，这样做虽然减少了监督频率，却不会自动得到好结果。原因在于标准 RL 在长时程、非回合式训练里会同时遇到两个问题。

\begin{figure}[H]
\centering
\includegraphics[width=\textwidth]{pdf_pages/page-17.png}
\caption{非回合式训练的两个核心挑战：探索让轨迹漂移离开任务分布，状态分布坍塌让 agent 学不到好策略。}
\end{figure}

第一，\textbf{exploration drift}。智能体在探索时很容易漂移到与任务无关、甚至难以恢复的状态区域。第二，\textbf{state distribution collapse}。当智能体离开有意义的状态簇之后，后续采样越来越集中在“坏状态”附近，更新也会围绕这些坏状态展开，最后甚至学不到有效策略。

\subsection{为何 episodic RL 能工作，而 reset-free RL 更难}

在标准 episodic setting 中，失败一次并不严重，因为下一次又会从初始状态分布 \(\rho_0\) 重新开始。这样训练数据天然围绕“可完成任务”的状态区域展开，学习信号始终比较集中。到了 non-episodic setting，这种隐式课程设计消失了。失败之后，agent 还得自己想办法回到有用状态，而这件事本身就是一个新的控制任务。

\subsection{先学会回去，才有机会继续学会向前}

讲者接着展示了一个经验现象：如果允许算法频繁回到初始状态，它能学得很好；一旦去掉频繁 reset，性能显著下降。这说明问题不在于 forward task 本身太难，而在于 agent 缺乏回到“可学习区域”的机制。

\begin{figure}[H]
\centering
\includegraphics[width=0.78\textwidth]{pdf_pages/page-18.png}
\caption{实验趋势说明：能够频繁回到初始分布时，forward task 能学会，因此关键难点转化为“如何学会恢复或重置”。}
\end{figure}

这正是本讲后续算法设计的出发点。比起一味提高探索强度，更有效的路线是显式学习恢复能力。

\subsection{本章小结}

单纯增加 horizon 不能解决真实世界自主训练问题。真正困难的是，agent 在失败后会漂移到陌生且难恢复的状态，从而让训练分布逐步劣化。因此，自主 RL 需要把“恢复”“回退”“重置”本身纳入学习对象。

\section{核心方法：学习重置、学习循环、学习恢复}

\subsection{Forward-Backward RL}

最核心的思路是：为任务同时学习一个前向策略 \(\pi_f\) 和一个后向策略 \(\pi_b\)。前向策略负责完成任务，后向策略负责回到接近初始分布的状态。两者交替 rollout、各自按对应奖励更新，从而减少对人工 reset 的依赖。

\begin{figure}[H]
\centering
\includegraphics[width=\textwidth]{pdf_pages/page-19.png}
\caption{Forward-Backward RL 的基本流程：前向策略做任务，后向策略学习回到初始区域；测试时只保留前向策略。}
\end{figure}

可以把它写成下面这个交替过程：
\[
\text{rollout }\pi_f \rightarrow \text{update }\pi_f \rightarrow \text{rollout }\pi_b \rightarrow \text{update }\pi_b.
\]

其中：
\begin{itemize}
\item \(\pi_f\) 的奖励 \(r_f(s,a)\) 鼓励到达目标区域；
\item \(\pi_b\) 的奖励 \(r_b(s,a)\) 鼓励回到起始状态分布附近；
\item 训练时尽量不 reset，只在必要时低频 reset；
\item 测试时丢弃 \(\pi_b\)，只部署 \(\pi_f\)。
\end{itemize}

这个方法的优点是概念简单，而且能直接把“恢复能力”做成一个可学习对象；缺点也很明确：后向策略在最终部署时会被丢弃，因此训练花费的一部分能力并不直接服务于最终任务。

\subsection{从“回到起点”到“回到专家状态”}

讲者随后介绍了一个重要变体：后向策略不一定非得回到单一初始状态，而可以回到示范轨迹或专家轨迹覆盖到的状态分布。这样做的动机是，真实任务往往存在一大片“有用但不完全等于起点”的恢复区域。只要能回到这些熟悉状态，forward policy 就可能继续完成任务。

这个修改本质上把 backward reward 从“reset to start”推广为“reset to good familiar states”，既减轻了恢复任务难度，也更适合用 demonstration 数据做辅助监督。

\subsection{任务循环：不只做前进和后退}

如果环境天然支持多个相互连接的子任务，那么更一般的做法不是一味学一前一后两种行为，而是构造任务循环，让 agent 在一个任务图中不断切换练习对象。

\begin{figure}[H]
\centering
\includegraphics[width=0.82\textwidth]{pdf_pages/page-29.png}
\caption{任务循环视角：把环境中的恢复、前进、切换都看成任务图上的边，从而避免把恢复限制成单一 backward task。}
\end{figure}

这比单纯的 forward-backward 更一般。比如开门任务可以和关门任务组成循环；拾起杯子、放回杯子、移动到新位置也可以组成任务链。训练时 agent 根据当前状态选择合适任务，而不是固定执行“做任务”或“撤回任务”。

\begin{importantbox}{从算法设计角度看}
本讲的共识是：真实世界自主训练需要显式建模可恢复性。恢复可以是 backward policy、expert-state reset policy，或者更一般的 task graph。
\end{importantbox}

\subsection{本章小结}

reset-free RL 的关键不是更激进地探索，而是学习怎样回到有用状态。Forward-Backward RL 提供了最直接的框架，而 expert-state reset 和 task-cycle 则把恢复从“回到起点”推广为“回到可继续学习的熟悉状态”。

\section{Single-Life RL：部署时也要在单次生命里适应}

\subsection{问题设定}

讲座最后把问题进一步推进到更困难的 single-life setting：智能体并不是从零开始训练，而是已经有过预训练经验；现在它进入新环境，只有“一条命”，要在单次部署过程中一边执行、一边适应，不能依赖 reset 反复试错。

这和前面的 autonomous RL 有联系，但更强调部署时在线适应。讲者用人开门失败后迅速调整钥匙位置的例子来说明，人类在部署中也会即时纠错，而不是回到训练室重新训练。

\subsection{为什么直接 test-time fine-tuning 不够}

一个简单基线是：部署时继续对原策略做 TD update 或 policy gradient update。讲者展示的现象是，这种做法有时能临时学会越过一个障碍，但一旦进入此前没见过的卡死状态，就可能再也爬不出来。原因在于：
\begin{itemize}
\item 在线更新仍然围绕当前目标奖励进行，未必鼓励先恢复到熟悉状态；
\item 低层动作空间中的微调幅度有限，很难实现策略层面的快速改道；
\item 进入坏状态后，局部数据分布太差，更新反而容易越调越偏。
\end{itemize}

\subsection{两个更有前景的方向}

讲者给出两类更有希望的思路。第一类，是在部署时把 RL 信号改成“回到熟悉状态”，而不是直接“继续冲向目标”。这样 agent 先把自己从坏状态拉回训练分布附近，再继续完成任务。第二类，是在更高层的技能空间里适应，而不是只在低层动作上微调。例如预训练“向前走”“翻身站起”“重新接近门把手”等技能，部署时在技能层面切换，会比逐动作修补更有效。

这类方法与现代 foundation model 也自然衔接：如果预训练模型本身带有较强的常识和技能先验，那么 single-life adaptation 可以少一些盲目试错，多一些结构化恢复。

\subsection{本章小结}

single-life RL 讨论的是更贴近真实部署的问题：系统已经学过一些东西，但在陌生条件下只有一次机会。此时仅靠低层在线微调通常不够，更关键的是恢复到熟悉状态以及调用更高层、可组合的技能。

\section{总结与延伸}

Lecture 16 的价值在于，它把真实机器人 RL 中最常被忽略的成本显式化了：\textbf{reset 不是免费的}。一旦接受这一点，问题就从“怎么学会一个好策略”扩展成“怎么在真实世界里持续地学、恢复、再继续学”。这也是为什么 autonomous RL 既关心最终策略质量，也关心生命期累计回报。

从方法上看，本讲给出了一条非常清晰的谱系：
\begin{itemize}
\item 最基础的是 Forward-Backward RL，用后向策略显式学习恢复。
\item 更实用的是把恢复目标放宽到示范覆盖的熟悉状态，而不是单一起点。
\item 更一般的是 task cycles，把恢复、重试和子任务切换统一进任务图。
\item 更前沿的是 single-life RL，要求系统在单次部署里继续适应。
\end{itemize}

如果把这讲和后续机器人课程线索连起来，可以得到一个很强的研究命题：真实世界 RL 的核心并不只是 sample efficiency，而是 \textbf{recoverability + adaptability}。未来更值得关注的方向包括如何自动发现恢复技能、如何用 demonstrations 或 world models 定义熟悉状态、以及如何把 LLM 或技能库用于单生命部署中的快速决策切换。

\end{document}
